\documentclass[12pt,english]{reviewresponse}

%% Language
\usepackage{babel}
\usepackage[babel]{microtype}
\usepackage[babel]{csquotes}
\usepackage{xcolor}

%% Fonts
\usepackage[T1]{fontenc}
\usepackage{lmodern}
%\usepackage{newcent} % different font
%\usepackage[scaled]{beramono} % different monospace font


%% Bibliography
\usepackage[backend=biber,style=ieee,dashed=false,url=false,isbn=false,defernumbers=true,refsection=section]{biblatex}
\bibliography{literature.bib}

\usepackage{hyperref}

\newcommand{\DC}[1]{{\color{violet}[{\tiny D:~}#1]}}
\newcommand{\CP}[1]{{\color{blue}[{\tiny CP:~}#1]}}


\title{Quantifying Very Extreme Precipitation and Temperature Using Huge Ensembles Generated by Machine Learning-based Climate Model Emulators}
\author{Christopher Paciorek and Daniel Cooley}
\journal{BAMS}

\begin{document}
\maketitle

\editor
\begin{generalcomment}
	On the basis of these reviews and my own evaluation, the manuscript requires major and mandatory revision. Acceptance is contingent upon satisfactorily addressing the concerns and suggestions of the reviewers.
\end{generalcomment}
\begin{revresponse}

Thank you for the opportunity to revise our manuscript. Below we address the comments of Reviewer 1 and we discuss the comments of Reviewer 2, in particular highlighting how our work builds on and goes beyond that of the Ben Alaya paper.

 


\end{revresponse}

\reviewer 

\begin{revcomment}
   I don't really understand the response to my Comment 3 in the original review—maybe I didn't explain myself well. What I mean is that, when you take a set of events, e.g., spacetime fields of precipitation associated with storms, and aggregate them to some spatiotemporal scale (e.g. 24-hour, watershed scale rainfall depth), fit an EV distribution, and then extrapolate to a particular AEP from that distribution, you are not able (at least not without additional nontrivial assumptions) to then "recover" spacetime information about that hypothetical event; all you have is, e.g., a 24-hour watershed-scale rainfall depth. Which is very useful, but is far from the complete story. For example, to run a flood simulation model (at least to run it well), you need high-resolution space-time rainfall fields. I have spent a lot of time arguing that this space-time information matters a lot and is a key reason to, if possible, prefer direct empirical AEP estimation over EV methods.
\end{revcomment}

\begin{revresponse}

%Thanks for this clarification. 
We agree the high-resolution fields are critical for many downstream applications including flood simulation. 
%For the scenario you present, 
For such applications, one could retain the space-time fields from the emulator for times (i.e., events) when precipitation is extreme 
%(or for all times, but isolating to extreme times would save computation) 
and use these as input to the flood simulation model to produce river flow values. 
One could then use EV methods to produce estimates of very low AEP events for flows instead of precipitation. 

However, we don't see why the need for full space-time fields is a reason to advocate for empirical AEP estimation rather than EV-based estimation. If the AEP estimate of interest is within the magnitude of values actually observed (e.g., the 1-in-1000 year AEP estimates if using emulator or model output), one could use the space-time fields from events whose overall magnitude is roughly that of the estimate, as you suggest. Even in this interpolation setting, EV methods could be better than empirical estimation, as using EV theory could reduce statistical variance. EV methods are required when one must extrapolate.
The amount of simulation/emulation required to produce multiple very extreme realizations in order to empirically estimate very low probability (e.g., 1-in-100K year) events (particularly under different climate scenarios (see comment 2)) would be massive.
It is beyond the scope of our work to discuss approaches one might use to produce reasonable space-time fields for very extreme events that correspond to extrapolation. 
But one might consider space-time EV-based modeling approaches to try to extrapolate the simulated extreme space-time fields to produce space-time fields that are even more extreme and which could be used for flood modeling. Or one might investigate some sort of pattern scaling from observed less extreme events.


In both the introduction and discussion, our previous submission already highlights the importance of the emulator producing full space-time fields (also emphasized in the NASEM report), so we have not modified the manuscript on this topic.

\end{revresponse}

\begin{revcomment}
   Lines 160-161: "Our main goal is to explore big-picture statistical questions, rather than to account for finer-grained sources of variability." That's understandable and ok here, but I do think there are some important questions about how emulators (or numerical models) should be run for these types of applications. For example, CESM2-LENS and other large ensemble climate model simulations that have a climate signal pose challenges related to internal and external (forced) variability when used to estimate nonstationary AEPs. We wrestled with this issue a bit, after from bruising peer review, in Liu et al. (2025). In contrast, the really interesting d4PDF ensemble simulations from Japan approach variability in a different way that is useful in these applications, by picking a particular radiative forcing level (e.g., corresponding to +2C global temperature rise) and running a big ensemble at that fixed level. In other words, if we're going to start running emulators or
numerical models for the purpose of AEP estimation, there needs to be active dialog between modelers and statisticians about the right way to force them, and then the right ways to analyze their outputs.
\end{revcomment}

\begin{revresponse}
Yes, we agree that the question of how to force the emulators is important. We touch on this from the perspective of accounting for nonstationarity. Our previous submission said "We see the main challenges as developing an appropriate emulator and deciding on the scenarios, not in the statistical analysis" (now in lines 398-399). Saying more is beyond the scope of our work, particularly in light of being up against the BAMS word limit.

\end{revresponse}

\begin{revcomment}
   L295: "by the time duration" is unclear. Usually in hydrometeorology we use duration to denote the length or averaging period of an event (e.g., annual max 24-hour rainfall, or a storm lasted 6 hours, etc.). But I think here you're referring to the simulation length (22 years) or the sample size? I'm not sure.
\end{revcomment}

\begin{revresponse}

Yes, it can be confusing because time comes in in various ways. Also in addition to line 295, we used "duration" in inconsistent ways elsewhere, so thanks for pointing this out. In line 295, we meant the return period/exceedance probability/time period of interest (1-in-10000 years, 1-in-100000 years, etc.). We now use "exceedance probability of interest" instead of "time duration" there (still at line 295). We also adjusted the language in several other places to avoid using "duration" (we now use "period" instead).

\end{revresponse}

\begin{revcomment}
   L307-308: it seems like the final thought here ("while cautioning that…") could benefit from a bit more description for the benefit of readers with less statistical prowess.
\end{revcomment}

\begin{revresponse}
We have reworded this in lines 308-309 to give some less technical framing: "while cautioning that the benefit of this diminishes because relative statistical uncertainty decreases based on the inverse of the square root of the sample size".

\end{revresponse}

\begin{revcomment}
   L372: "functional"? I'm not sure what you mean from the context.
\end{revcomment}

\begin{revresponse}

The use of "functional" here is statistical terminology that refers to a function of a distribution. We've changed to "function" to be more clear (line 373).

\end{revresponse}

\reviewer 
\begin{revcomment}
I appreciate the authors' efforts in revising their manuscript. It is now clear that the focus of the paper is to demonstrate the use of large ensembles to estimate very extreme statistics. The authors show that when a sufficiently high threshold is applied in threshold‑exceedance methods, extremely rare values can be estimated reliably. They also find that their conclusions are consistent with the earlier work of Ben Alaya (2020), which used the block‑maximum method.

However, because this is the central contribution of the paper, the advancement appears quite incremental. The authors have simply conducted a analysis similar to Ben Alaya (2020) with a different fitting technique on a different dataset, and the findings are of course entirely expected. Given this, I do not see sufficient novelty to warrant publication in BAMS.
   
\end{revcomment}

\begin{revresponse}

We agree that our work has some similarity to Ben Alaya et al. (2020) (henceforth "BA") and provides additional support/replication for its conclusions, which focus on statistical bias. However, our work goes beyond BA in the following ways:

\begin{itemize}
\item Our results are based on a machine learning  emulator trained on ERA5 reanalysis rather than traditional climate modeling as in BA. Such emulators are the subject of great current interest and have high potential for use in the future, so this is a particularly relevant "different dataset". One of our results is that the ML emulator is able to produce out-of-training-sample extremes, although our focus is not on a full assessment of fitness-for-purpose.
\item Our work considers the statistical questions in the context of an important real-world application, that of PMP, the subject of the recent NASEM report and ongoing work by NOAA and partners to update PMP estimates, as required by US federal legislation.
\item We systematically address the practical question of how to deal with seasonality/storm types. From first-hand experience during the NASEM report work, we know that uncertainty about how to address this in a principled way is common.
\item We discuss sample size considerations.
\item Our sample size of 10560 years is nearly an order of magnitude larger than the 1750 years in BA, providing for more robust analysis when using very high thresholds (or, equivalently, long blocks) and allowing us to compare EVA-based 1-in-10000 year AEP depths to direct empirical estimates, whereas the BA work goes only as far as 1-in-1000 year depths. For the PMP application, the difference between 1000 and 10000 years is material; indeed, PMP applications may need 1-in-100000 or 1-in-million year estimates.  
\item We report an intriguing result for precipitation regarding systematic decreases in the estimated shape parameter with increasing thresholds that, while preliminary, raises interesting questions about the boundedness of extreme precipitation.
\item We include consideration of temperature, which has fundamentally different behavior than precipitation.
\end{itemize}

In summary,  while the BA work is an important forerunner to our work, our work is not simply incremental. It provides important additional results that go beyond those reported in BA and are useful for real applications, which are likely to use emulator output in the future. 

We recognize this response may not fully satisfy reviewer 2's concerns, but we do not 
see a path to further modifying the submission to do so.
Given that the two other reviewers support publication and given the original interest from the BAMS editorial staff in the topic as presented in our BAMS paper proposal, we believe the work would be of interest to the BAMS readership.



\end{revresponse}

\reviewer 
\begin{revcomment}
I have examined the responses by the authors and have also cross-checked key passages in the text.  I am satisfied with the revisions.

\end{revcomment}

\begin{revresponse}
No response needed.
\end{revresponse}

\end{document}